\section{Three-layer benchmark}

\begin{table}[t]
\centering
\caption{Pre-registered benchmark layers. No dataset or result count is inferred here.}
\label{tab:datasets}
\small
\begin{tabular}{llll}
\toprule
Layer & Family & Shift surface & Current status \\
\midrule
Fraud & Elliptic & time / graph access / budget & saved outputs verified; pilot pending \\
Fraud & DGraphFin & time / graph access / resource & saved outputs verified; pilot pending \\
Controlled & fifteen mechanisms & isolated and correlated factors & tiny fixtures ready \\
Graph OOD & GOOD & official environments & not installed \\
Fallback & OGB molecular & scaffold / time & licence review pending \\
\bottomrule
\end{tabular}
\end{table}

\begin{table}[t]
\centering
\caption{Canonical saved-output expert set for the first fraud pilot.}
\label{tab:experts}
\small
\begin{tabular}{llll}
\toprule
Expert & Information & Resource sensitivity & Evidence status \\
\midrule
Feature MLP & node features & low graph dependence & members verified; result pending \\
GCN & features + visible graph & graph / memory / latency & members verified; result pending \\
GraphSAGE & features + sampled graph & graph / sampling / latency & members verified; result pending \\
\bottomrule
\end{tabular}
\end{table}


\paragraph{Layer A: fraud deployment shifts.}
The first saved-output pilot uses Elliptic and DGraphFin under strict inductive, isolated inductive, and transductive-structure protocols. Each protocol is mapped to explicit time, visibility, and construction axes rather than treated as an opaque name. For each dataset, target protocol, and seed, the target protocol is held out and the remaining two provide six source expert bindings; the target provides three test bindings. The resulting 180 role-neutral coordinates, 60 scenarios, and 540 bindings are registered. All six canonical archives and 180 members pass checksum, schema, coordinate, chronology, label-known, cross-expert alignment, and cross-protocol row-scope audits; the pilot itself has not run.

Future fraud conditions add temporal drift, graph construction/visibility changes, expert unavailability, review-budget contraction, latency/memory limits, and combined shifts. Provider task units and label mappings are retained. Unknown labels are excluded; chronological and endpoint scope must validate before scoring.

\begin{figure}[t]
  \centering
  \includegraphics[width=\linewidth]{05_scenario_construction.pdf}
  \caption{Scenario-local roles. A base artifact may change role across scenarios but cannot be both source and target inside one scenario.}
  \label{fig:scenario}
\end{figure}

\paragraph{Layer B: controlled mechanisms.}
Fifteen deterministic generators isolate feature corruption, edge deletion, degree shift, homophily reduction, heterophily reversal, temporal concept drift, delayed labels, neighborhood truncation, expert unavailability, review-budget contraction, correlated graph/feature shift, noisy metadata, missing axes, dynamic resources, and universal expert degradation. Every mechanism declares source and held-out contracts, qualitative router direction, pilot/full sizes, seeds, metrics, failure modes, and deterministic hash. Tiny fixtures test semantics; they do not claim fraud realism.

\begin{figure}[t]
  \centering
  \includegraphics[width=\linewidth]{08_synthetic_mechanism_taxonomy.pdf}
  \caption{Controlled mechanism taxonomy. Expected directions are hypotheses, not hard-coded empirical conclusions.}
  \label{fig:mechanisms}
\end{figure}

\paragraph{Layer C: non-fraud graph OOD.}
GOOD is the primary recognized family because it exposes official graph OOD environments and task-compatible node settings \cite{gui2022good}. An OGB molecular scaffold/time family is a fallback subject to dataset-specific licence review and a graph-level expert bridge. No official code or dataset is installed in this build. Both adapters have tiny stand-ins that test schema, label access, and split handling but cannot support an official result claim.

\paragraph{Interventions and leakage controls.}
The benchmark includes held-out combinations, unseen values, noisy/partial metadata, counterfactual contracts, resource and budget masks, all-expert failure, and all-unavailable cases. Contract changes never mutate fixed expert scores during counterfactual routing tests. Dataset, protocol, expert, seed, fold, row ID, provider labels, scores, split, and label-known status must align before evaluation.
